\documentclass[a4paper, 12pt, oneside, headspline]{scrartcl}
% !TeX root = ../main.tex
% Variablen welche innerhalb der gesamten Arbeit zur Verfügung stehen sollen
\newcommand{\titleDocument}{Bachelorarbeit}
\newcommand{\subjectDocument}{Test-driven development mit Python}
\newcommand{\institution}{Hochschule für angewandte Wissenschaften Augsburg}
\newcommand{\insSection}{Fakultät für Informatik}
\newcommand{\name}{Maximilian Konter}
\newcommand{\email}{maximilian.konter@hs-augsburg.de}

% Deaktiviere den indent nach jedem Absatz
\setlength{\parindent}{0pt}

\usepackage[openbib]{currvita}

% wichtig für Trennung von Wörtern mit Umlauten
\usepackage[T1]{fontenc}

% verbesserter Randausgleich
\usepackage{microtype}

% Farbige links
\usepackage{color}

% Umlaute unter UTF8 nutzen
\usepackage[utf8]{inputenc}

% deutsche Silbentrennung
\usepackage[ngerman]{babel}

% Grafiken aus PNG Dateien einbinden
\usepackage{graphicx}

% bricht lange URLs "schoend" um
\usepackage[hyphens,obeyspaces,spaces]{url}

% mehrseitige Tabellen ermöglichen
\usepackage{longtable}

% für Tabellen
\usepackage{array}

% Schaltet den zusätzlichen Zwischenraum ab, den LaTeX normalerweise nach einem 
%Satzzeichen einfügt.
\frenchspacing

% Paket für Zeilenabstand
\usepackage{setspace}

% für Bildbezeichner
\usepackage{capt-of}

% für Stichwortverzeichnis
\usepackage{makeidx}

% Für Glossar
\usepackage{glossaries}

% Ermöglicht das Fußnotenverzeichnis
\usepackage[titles]{tocloft}

%% Fußnotenverzeichnis von https://stackoverflow.com/a/6914775/8339179
\newcommand{\listfootnotesname}{Fußnotenverzeichnis}%     % 'List of Footnotes' 
%%title 
\newlistof[chapter]{footnotes}{fnt}{\listfootnotesname}%  % New 'List of...' 
%%for footnotes 
\let\oldfootnote\footnote%                                % Save the old 
%%\footnote{...} command 
\renewcommand\footnote[1]{%                               % Redefine the new 
%%footnote to also add 'List of Footnote' entries. 
    \refstepcounter{footnotes}%                           % Add and step a 
    %%reference to the footnote/counter. 
    \oldfootnote{#1}%                                     % Make a regular 
    %%footnote. 
    \addcontentsline{fnt}{footnotes}{\protect% 
        \numberline{\thefootnotes}#1}%                    % Add the 'List 
        %%of...' entry. 
}

% Ermöglicht mehr als 3 punkte tiefe -> 1.2.3.4
\setcounter{secnumdepth}{4}
\setcounter{tocdepth}{4}

\newcommand{\myparagraph}[1]{\paragraph{#1}\mbox{}\\\noindent}

% Von https://tex.stackexchange.com/a/10110 mit additions aus 
%https://stackoverflow.com/a/6914775/8339179
\newcommand{\footlabel}[2]{%
    \refstepcounter{footnotes}%
    \addcontentsline{fnt}{footnotes}{\protect 
        \numberline{\thefootnotes}#2}%
    \addtocounter{footnote}{1}%
    \footnotetext[\thefootnote]{%
        \addtocounter{footnote}{-1}%
        \refstepcounter{footnote}\label{#1}
        #2%
    }%
    $^{\ref{#1}}$%
}

\renewcommand{\footref}[1]{%
    $^{\ref{#1}}$%
}


% Bessere Listen
\usepackage{enumitem}

\usepackage[
    maxnames=6,
    minnames=1,
    style=apa,
    backend=biber,
    bibencoding=utf8,
    sortlocale=de_DE,
]{biblatex}
\DeclareLanguageMapping{ngerman}{ngerman-apa}
\DefineBibliographyStrings{ngerman}{
    andothers = {{et\,al\adddot}},
}

% für Listings
\usepackage{listings}
\lstset{
    numbers=left,
    numberstyle=\tiny,
    numbersep=5pt,
    keywordstyle=\color{black}\bfseries,
    stringstyle=\ttfamily,
    showstringspaces=false,
    basicstyle=\footnotesize,
    captionpos=b
}

\newglossaryentry{Pythonista}{
    name={Pythonista},
    description={Jemand, der Python als seine favorisierte Sprache verwendet}
}
\newglossaryentry{Continous testing}{
    name={Continous testing},
    description={Bezeichnet den Prozess des automatischen Testens mit Hilfe
    eines Automatisierungstools. Hierbei werden die Tests durch ein Ereignis
    automatisch ausgeführt, wodurch der Entwickler sofort ein Feedback bekommt}
}
\newglossaryentry{mock}{
    name={mock},
    description={Etwas mocken bedeutet, ein Objekt durch ein falsches Objekt, den
    Mock zu ersetzen, das genau so aussieht, wie das erwartete Objekt, aber nicht das
    Gleiche ist. So ist es möglich eine Funktion zu mocken, wodurch nicht die Funktion,
    sondern der Mock aufgerufen wird und das eingestellte ergebnis liefert}
}
\newglossaryentry{commit}{
    name={commit},
    description={Ein commit ist die Sammlung von änderungen an Dateien, welche 
    mithilfe eines VCS verwaltet werden},
    plural=commits
}
\newglossaryentry{vcs}{
    name={Versions Control System},
    description={Ein VCS ist ein Dienst, der es seinen Nutzern ermöglicht änderungen
    an Dateien in einer Chronik zu speichern um später darauf zu zu greifen. Dies
    dient auch der verteilung von Daten über mehrere Systeme sowie sicherung der
    Daten vor verlust}
}
\newglossaryentry{fuzz}{
    name={fuzz-testing},
    description={Beschreibt eine Art von Test,
    bei dem das zu testende Modul oder Programm mit zufälligen Werten aufgerufen wird. Dabei
    soll jede mögliche anwendung des Moduls oder Programms dargestellt werden}
}
\newglossaryentry{bug}{
    name={bug},
    description={Ein Bug(zu deutsch Käfer) ist ein Fehler in einem Programm oder in einer Software}
}
\newglossaryentry{fixture}{
    name={fixture},
    description={Beschreibt die Präperation von Tests mit festen (fix) Werten, so wird zum Beispiel
    vor jedem Test eine Variable gesetzt. Dies dient der übersichtlichkeit, da dies dann nicht in
    jedem Test vorgenommen werden muss und um varianz zwischen Tests und Test läufen zu vermeiden}
}
\newglossaryentry{cli}{
    name={Command-line interface},
    description={Eine Benutzerschnitstelle für das Terminal}
}
\newglossaryentry{stacktrace}{
    name={StackTrace},
    description={Der StackTrace ist ein Protokoll der aufgerufenen Funktionen und Methoden die 
    ineinander verschachtelt sind, bis hin zu der tiefsten Funkion/Methode die einen Fehler
    geworfen hat}
}

% Abkürzungsverzeichnis
\usepackage[german]{nomencl}
\let\abbrev\nomenclature

% Abkürzungsverzeichnis LiveTex Version
\renewcommand{\nomname}{Abkürzungsverzeichnis}
\setlength{\nomlabelwidth}{.25\hsize}
\renewcommand{\nomlabel}[1]{#1 \dotfill}
\setlength{\nomitemsep}{-\parsep}
\makenomenclature
\makeglossaries

% Disable single lines at the start of a paragraph (Schusterjungen)
\clubpenalty = 10000
% Disable single lines at the end of a paragraph (Hurenkinder)
\widowpenalty = 10000
\displaywidowpenalty = 10000

% Hyperlinks im Text
\usepackage[bookmarksnumbered,pdftitle={\titleDocument},hyperfootnotes=false]{hyperref}
%\hypersetup{colorlinks, citecolor=red, linkcolor=blue, urlcolor=blue}

% neue Kopfzeilen mit fancypaket
\usepackage{fancyhdr}                  % Paket laden
\pagestyle{fancy}                      % eigener Seitenstil
\fancyhf{}                             % alle Kopf- und Fußzeilenfelder 
%bereinigen
\fancyhead[L]{\nouppercase{\leftmark}} % Kopfzeile links
\fancyhead[C]{}                        % zentrierte Kopfzeile
\fancyhead[R]{\thepage}                % Kopfzeile rechts
\renewcommand{\headrulewidth}{0.4pt}   % obere Trennlinie
\fancyfoot[C]{\thepage}                % Seitennummer
\renewcommand{\footrulewidth}{0.4pt}   % untere Trennlinie

% Bibliothek hinzufuegen
\addbibresource{bachelor.bib}

\begin{document}
    % Titelseite
    % !TeX root = ../bachelor.tex
\thispagestyle{empty}

\begin{figure}[t]
 \centering
 \includegraphics[width=0.3\textwidth]{abb/logoHSA}
\end{figure}


\begin{verbatim}


\end{verbatim}

\begin{center}
\Large{\institution}
\end{center}


\begin{center}
\Large{\insSection}
\end{center}
\begin{verbatim}




\end{verbatim}
\begin{center}
\doublespacing
\textbf{\LARGE{\titleDocument}}
\singlespacing
\begin{verbatim}

\end{verbatim}
\textbf{{~\subjectDocument~}}
\end{center}
\begin{verbatim}

\end{verbatim}
\begin{center}

\end{center}
\begin{verbatim}

\end{verbatim}
\begin{center}
\textbf{zur Erlangung des akademischen Grades \\ Bachelor of Science}
\end{center}
\begin{verbatim}





\end{verbatim}

\begin{footnotesize}
\begin{flushleft}
\begin{tabular}{llll}
\textbf{Thema:} & & \subjectDocument & \\
& & \\
\textbf{Autor:} & & \name & \\
& & \email & \\
& & MatNr. 951004 & \\
& & \\
\textbf{Version vom:} & & \today &\\
& & \\
\textbf{1. Betreuer:} & & Dipl.-Inf. (FH), Dipl.-De Erich Seifert, MA &\\
\textbf{2. BetreuerIn:} & & Prof. Dr.-Ing. Alexandra Teynor &\\
\end{tabular}
\end{flushleft}
\end{footnotesize}

    % 1.5 facher Zeilenabstand
    \onehalfspacing

    % Einleitung / Abstract
    \renewcommand{\abstractname}{Abstract}
    % !TeX root = ../bachelor.tex
\begin{center}
    \vspace{0.9cm}
    \textbf{Abstract}
\end{center}

\begin{abstract}
Da die Anforderungen an Software im Laufe der Zeit stark angestiegen sind,
haben sich einige Technologien entwickelt, die dazu verwendet werden diese
Anforderungen zu gewährleisten. Eine dieser Technologien ist das Test-driven
development. Diese Arbeit befasst sich mit dieser Technologie im Bezug zu den
aktuellen Test-Tools der Programmiersprache Python und analysiert diese auf
ihre Tauglichkeit im Bezug zum Test-driven development.
\newline

Dabei werden zuerst verschiedenen Arten von Tools untersucht, die unit-testing-,
mock-testing- und fuzz-testing Tools. Diese werden anhand ihrer
Anwendbarkeit, Effizient, Komplexität und Erweiterbarkeit beschrieben und
analysiert, um mit Hilfe eines Vergleichs derer einen Leitfaden für das am
besten für Test-driven development geeignete Tool zu bieten. Das daraus
resultierende Ergebnis soll aufzeigen, welche Toolkombinationen die beste Wahl
anhand der vorgegebenen Kriterien darstellen. Diese sind Anwendbarkeit,
Effizienz, Komplexität und Erweiterbarkeit.
\newline

Die dabei entstandenen Ergebnisse zeigen auf, dass nicht alle Tools optimal
für Test-driven development geeignet sind. Schlussendlich ergaben sich vier
Tools, welche gemeinsam eingesetzt, das optimale Toolset für die allgemeine
Anwendung bieten. Dabei ist jedoch zu beachten, dass dies immer vom
Anwendungsfall, sowie den Anforderungen an die Software und an die Tools
abhängig ist. Dabei ergab sich, dass \lstinline{pytest} zusammen mit
\lstinline{mocktest} sowie \lstinline{doctest} und \lstinline{hypothesis}
gemeinsam eine sehr gute Wahl darstellen.
\end{abstract}


    % einfacher Zeilenabstand
    \singlespacing

    % Inhaltsverzeichnis anzeigen
    \newpage
    \tableofcontents

    %%%%% Beginn Dokument %%%%%
    \newpage
    \fancyhead[L]{\nouppercase{\leftmark}} % Kopfzeile links

    % 1,5 facher Zeilenabstand
    \onehalfspacing

    % einzelne Kapitel
    % !TeX root = ../bachelor.tex
\section*{Einleitung}\label{einleitung}
\addcontentsline{toc}{section}{Einleitung}

Die Ansprüche an Software sind in den letzten Jahren immer weiter gestiegen.
Dies liegt vor allem an der Reichweite, die Software heute hat. So besitzen
im Jahr 2018 bereits circa 66\% aller Menschen ein Smartphone
(\cite{FraukeSchobelt:Smartphone}) und im Arbeitsleben ist ein PC meist nicht
mehr weg zu denken. Die dadurch benötigte Software muss bestimmte
Anforderungen, welche die Nutzer an diese stellen, gewährleisten und
sicherstellen. Mit steigenden Nutzerzahlen, steigen oft auch die Anforderungen
an eine Software, sowie die gefundenen \Glspl{bug} in dieser.

Im weltweiten Markt gibt es viele große Unternehmen, die gegenseitig um ihre
Nutzer kämpfen. Die Anwender präferieren meist denjenigen Anbieter, welcher die
bessere Software bietet. An dieser Stelle ist der Software-Architekt einer
jeden Anwendung gefragt, welche Technologie eingesetzt wird um die qualitativen
Anforderungen an diese zu gewährleisten. Eine dieser Technologien ist das
Test-driven development, welches sich durch eine hohe resultierende
Test-Abdeckung ihren Namen gemacht hat. Bei dieser Methode werden zuerst die
Tests und dann der Code geschrieben, um so später sicher stellen zu können,
dass die Software alle Anforderungen erfüllt und keine \Glspl{bug} enthält.

Gerade zum Release einer neuen Software ist es für Unternehmen schwer, Nutzer
zu akquirieren. Andere Anwendungen, die bereits auf dem Markt sind, haben
schon viele ihrer \Glspl{bug} gefixed und verfügen über eine durch den
Anwender ausgiebig getestete Software. Um eine neue Applikation auf diesen
Markt zu bringen, ist es erforderlich, dass diese zum Start ohne Fehler läuft.
Dies kann mit Test-driven development sicher gestellt werden.

Eine der Programmiersprachen mit der Test-driven development möglich ist, ist
Python. Python wird heutzutage von vielen Entwicklern immer häufiger eingesetzt
und ist Mittlerweile die am zweit meisten beliebte und am vierthäufigsten
genutzte Sprache weltweit (\cite{stackoverflow:2019}).

Aufgrund der zunehmenden Beliebtheit, sowie meiner eigenen Vorliebe für die
Programmiersprache Python (auch \Gls{Pythonista} genannt) habe ich
mich dazu entschieden meine Bachelorarbeit dieser zu widmen. Da ich beginnend
mit dem Praxissemester und danach als Werkstudent in einer Firma an der
Entwicklung der hausinternen Test-Umgebung mitgewirkt habe, wurde ich immer
interessierter an dem Thema Testen. Durch die Beschäftigung mit Tests in meiner
Freizeit bin ich auf das Test-driven development gestoßen und war
sofort überzeugt, dass dies eine Technologie ist, die ich anwenden möchte. Doch
leider war die Auswahl der Tools, welche für Test-driven development in Frage
kamen so hoch, das ich Angst hatte etwas verpassen zu können, da ich ein Tool
einem anderen vorzog. Beim Analysieren der Tools viel mir auf, dass viele der
Tools veraltet oder ungeeignet waren, was einen Einsatz für Test-driven
development ausschloss. Die übrig gebliebenen Tools waren so umfangreich, dass
ich das Projekt erst einmal zur Seite legte.

Die Bachelorarbeit erschien mir als der passende Rahmen für die intellektuelle
und praktische Beschäftigung mit dieser Thematik. Zum einen konnte ich so meine
Privaten Pläne diesbezüglich verwirklichen und zum anderen erschien mir das
Thema relevant für meine späteren Beruflichen Perspektiven.

Das Ziel dieser Arbeit ist es, verschiedene aktuelle Test-Tools der
Programmiersprache Python auf ihre Tauglichkeit für den Einsatz von Test-driven
development zu untersuchen und zu vergleichen. Dabei wird zuerst jedes Tool und
seine Funktionalitäten beschrieben. Im Anschluss wird ein jedes auf die
verschiedene Anforderungen untersucht, welche es erfüllen muss um für
Test-driven development in Frage zu kommen. Dies geschieht anhand eines
erstellten Codebeispiels, welches gewisse Funktionalität liefern sollte.
Dieses Beispiel wird mithilfe der verschiedenen Tools auf die gleiche Art
getestet. Die dabei entstehenden Komplikationen, sowie der benötigte Aufwand
zeigen auf, welche Tools für die Analyse benötigt werden.

Die Arbeit soll aufzeigen welche Kombination der verschiedenen Tools die
optimale Verbindung bieten um Test-driven development durch zu führen. Dabei
erwarte ich mir tiefere Einblicke in das Test-driven development zu erhalten,
um heraus zu finden ob diese Technologie für mich und für meine weiteren
Arbeiten relevant ist.

Dazu wird in Kapitel \ref{definition} zunächst geklärt, welche Eigenschaften
die Programmiersprache Python hat, sowie die Definition von Test-driven
development. Die genutzte Methodik wird in Kapitel \ref{methodik} beschrieben
und anschließend in Kapitel \ref{python-tools} angewendet. Die Analyse ist in
vier Unterpunkte gegliedert, in Punkt \ref{python-tools:stdlib} werden zunächst
die Tools der Standardbibliothek analysiert. Diese sind \lstinline{unittest}
und \lstinline{doctest}. Zusammengefasst unter Punkt \ref{python-tools:extlib}
sind die Tools, welche nicht in der Standardbibliothek zu finden sind. Diese
setzen sich aus \lstinline{pytest}, \lstinline{stubble}, \lstinline{mocktest},
\lstinline{flexmock}, \lstinline{doublex} und \lstinline{hypothesis}.
Darauffolgend werden unter Punkt \ref{python-tools:vergleich} die analysierten
Tool verglichen um unter Punkt \ref{python-tools:kombination} mögliche
Kombinationen dieser auf zu Zeigen. Bevor in Kapitel \ref{fazit} das Fazit
gezogen wird, werden die Ergebnisse der Arbeit in Kapitel \ref{diskussion}
diskutiert.
    % !TeX root = ../bachelor.tex
\section{Definition}\label{definition}
% !TeX root = ../bachelor.tex
\subsection{Die Programmiersprache Python}\label{einleitung:python}

Python ist eine dynamisch typisierte Programmiersprache, sie ist Open Source
unter der Python Software Foundation License
(PSFL\footnote{https://docs.python.org/3/license.html}). Eine dynamisch
typisierte Programmiersprache bestimmt zur Laufzeit, welchem Typ eine Variable
angehört. Dies ist möglich, da Python eine interpretierte und keine kompilierte
Sprache ist. Interpretierte Sprachen werden, zur Laufzeit, in Maschinencode
umgewandelt, während kompilierte Sprachen davor umgewandelt werden. Beides hat
seine Vor- und Nachteile auf diese hier aber nicht weiter eingegangen werden
sollen, da dies nicht die Aufgabe dieser Arbeit ist.

Python unterstützt verschiedene Programmier-Paradigmen, wie die funktionale
Programmierung oder die objektorientierte Programmierung. Dadurch, dass Python
interpretiert ist, wird die Sprache auch oft als Skript Sprache genutzt.

Das Hauptmerkmal der Sprache ist die Syntax, die dabei verwendet werden muss.
Bei Python werden keine geschweiften Klammern genutzt, stattdessen werden
Einrückungen in Form von vier Leerzeichen verwendet. Dadurch entsteht ein
leicht zu lesender Code, der gerade für Anfänger besser zu verstehen ist.
% !TeX root = ../bachelor.tex
\subsection{Test-driven development}\label{einleitung:tdd}
Das Test-driven development (TDD) wurde durch Kent Beck, einem Amerikanischen 
Softwareingenieur und Autor bekannt. Dieser beschreibt in seinem Buch "`Test 
Driven Development: By Example"' wie bei der Anwendung von TDD vorgegangen 
werden soll. Demnach sieht der Entwicklungsprozess nach \cite{beck:tdd} 
folgendermaßen aus (siehe Abbildung \ref{abb:einleitung:tdd:zyklus}):
\newline

\begin{figure}
\includegraphics[width=1\textwidth]{abb/tdd.png}
\caption[By Xarawn - Own work, CC BY-SA 4.0,
\url{https://commons.wikimedia.org/w/index.php?curid=44782343}]{TDD Zyklus}
\centering
\label{abb:einleitung:tdd:zyklus}
\end{figure}

Am Anfang sollte geklärt sein, welche Funktionen erfüllt werden müssen. Anhand
dieser beginnt dann die Anwendung von TDD. Dazu wird der folgende Zyklus
beschrieben, der sich so lange wiederholt, bis alle Features implementiert
und getestet wurden.

Der Erste Schritt (siehe Abbildung \ref{abb:einleitung:tdd:zyklus} \#1) dazu ist
das Schreiben eines Tests, der neue Funktionalität oder Verbesserung zu einer
bestehenden Unit hinzufügt. Dieser Test sollte möglichst kurz und
aussagekräftig sein. Dazu muss der Entwickler die Spezifikationen und
Anforderungen des Features genau verstehen um den Test wirklich effektiv
schreiben zu können. Ein Entwickler, der die Anforderung nicht ganz versteht,
ist nicht in der Lage einen Test zu schreiben, welcher alle Bereiche abdeckt.
Dieser Entwickler muss allerdings nicht zwangsläufig später auch die
Funktionalität implementieren. Der Test sollte demnach so umfassend und genau
sein, dass jeder Entwickler die dazugehörige Logik implementieren kann.

Als nächstes (siehe Abbildung \ref{abb:einleitung:tdd:zyklus} \#2) werden alle
Tests die bereits bestehen, sowie der Neue durchgeführt um zu sehen, ob der
neue Test fehlschlägt. Dies bestätigt unter anderem auch, dass der neue Test
die anderen Tests nicht beeinflusst und auch ohne neuen Code, der
Funktionalität implementiert, nicht bestanden wurde. Dadurch kann der
Entwickler ausschließen, dass ein Test immer korrekt ist, wodurch das Vertrauen
des Entwicklers in den Test gesteigert wird.

Der nächste Schritt (siehe Abbildung \ref{abb:einleitung:tdd:zyklus} \#3) ist
die Implementierung des Codes, der getestet werden soll. Dabei ist unwichtig,
wie der Code geschrieben wird. Das bedeutet, dass es erlaubt ist
"`schmutzigen"' (Code der nicht den gefordertem Standard entspricht) Code zu
produzieren, welcher so nicht akzeptiert werden würde. Die Hauptsache dabei
ist, dass der Test bestanden wird. Die Codequalität spielt hier keine Rolle, da
sie in Punkt Fünf überarbeitet wird. Der Entwickler darf in diesem Prozess auch
keinen weiteren Code, der nicht notwendig ist um den Test zu bestehen,
hinzufügen. Dies verhindert, dass Code geschrieben wird, der keine zugehörigen
Tests hat.

Ist die gesamte Logik implementiert, werden im vorletzten Schritt (siehe 
Abbildung \ref{abb:einleitung:tdd:zyklus} \#4) alle Tests durchgeführt
und überprüft, ob diese erfolgreich waren. Sollte dies nicht der Fall sein,
muss die Implementierung korrigiert werden und/oder eventuell der Test 
angepasst werden (siehe Abbildung \ref{abb:einleitung:tdd:zyklus} \#5').

Sind alle Tests bestanden wird im finalen Schritt (siehe Abbildung 
\ref{abb:einleitung:tdd:zyklus} \#5) der neue Code aufgeräumt.
Dieser Punkt ist sehr wichtig, da in Punkt Drei der Code nur funktionieren 
muss. Dadurch kann es sein, dass hier Code entstanden ist, der nicht den 
qualitativen Anforderungen entspricht. Sämtliche Objekte sollten einen 
aussagekräftigen Namen erhalten und der Code sollte, sofern dies nötig ist an 
einen Ort verlegt werden, der seiner logischen Aufgabe entspricht. 
Duplikationen müssen entfernt werden und nach jeder Aufräumaktion sollten die 
Tests noch einmal ausgeführt durchlaufen, um zu verifizieren, dass alles noch 
funktioniert.
\newline

Dieser Zyklus wird nun von vorne so lange wiederholt, bis jedes Feature 
implementiert ist und die Software als fertig gestellt gilt. Die dafür nötigen 
Schritte sollten möglichst klein sein, um eine höhere Testabdeckung zu 
erreichen.
    \section{Methodik}\label{methodik}
Die zu analysierenden Tools, werden zwischen STDLIB und externen Tools 
unterschieden. Zusätzlich findet eine Gliederung in unit-,\gls{mock}- und 
\gls{fuzz} Tools statt.

Die unit-testing Tools sind Tools, die Funktionalität zum Testen bereitstellen.
Jedoch wird für TDD weit mehr als nur Tests benötigt. Aus diesem Grund werden
\gls{mock}-testing - und \gls{fuzz} Tools zusätzlich behandelt. Dabei soll
zwischen einer reinen Erweiterung eines Tools und der Erweiterung von allen
Tools unterschieden werden. Ist zum Beispiel ein Tool nur zusammen mit einem
anderen Tool, welches die Funktionalität bereit stellt Tests aus zu führen,
so ist dieses Tool als Erweiterung zu sehen und wird nicht analysiert. Bietet
ein Tool allerdings Funktionen zum Erweitern von verschiedenen Test Tools, so
wird es hier zur  Analyse verwendet. Diese Unterscheidung wird genutzt um
Duplikationen in den einzelnen Vergleichen zu vermeiden und verhindert, dass
Tools mit sehr vielen Erweiterungen den Rahmen dieser Arbeit sprengen.Sollten
sich für ein Tool besonders interessante Erweiterungen finden, so werden diese
in der Analyse des jeweiligen Tools erwähnt und verlinkt.
\newline

Jedes Tool wird anhand folgender Aspekte untersucht. Anwendbarkeit: Bietet das
Tool alles, um TDD betreiben zu können? Bei unit-testing Tools beinhaltet dies
\Glspl{fixture} und \Glspl{mock} während es bei den \gls{mock}- und \gls{fuzz} 
Tools um die verschiedenen Features geht. Sowohl das \Gls{mock}en als auch das
\Gls{stub}en sind für das TDD von großer Bedeutung. Einerseits ist es wichtig
nach zu verfolgen, welches Modul wann und mit welchen Parametern aufgerufen
wurde. Andererseits ist es genau so wichtig Tests ab zu schotten. Dies
geschieht mithilfe der \Glspl{stub}, welche es ermöglichen externe
Abhängigkeiten zu ersetzten um Fehler aus anderen Quellen aus zu schließen 
(\cite{percival:tdd:python}).
Zusätzlich wird in diesem Aspekt auch die benötigten Abhängigkeiten eines
Pakets untersucht, um dieses betreiben zu können. Mehr
Abhängigkeiten führen dazu, dass mehr Entwicklern vertraut werden muss, dass
diese Ihre Pakete aktuell und fehlerfrei halten.
Der zweite Aspekt ist die Effizienz mit der sich ein Tool einsetzen lässt. Dabei
wird darauf geachtet, wie viel Aufwand benötigt wird, das Tool ein zu setzen. Im
Besonderen wird hier auf die Vorarbeit, die benötigt wird um das Tool verwenden
zu können, geachtet. Des weiteren wird in diesem Aspekt bei den unit-testing
Tools die Effizienz, mit der Tests ausgewertet werden können, analysiert.
Als dritter Aspekt wird die Komplexität des Tools bewertet. Dabei wird diese
nicht zwangsläufig als negativ gesehen. Hat ein Tool viel
Funktionalität neben seiner Basis Features, so hat es in diesem Bezug eine hohe
Komplexität, die positiv zu sehen ist. Hingegen wird viel oder
unübersichtlicher Code den das Tool produziert als negativer Komplex gesehen.
Zuletzt wird bei den unit-testing Tools auf die Erweiterbarkeit geachtet. Diese
umfasst Erweiterungen der Entwickler als auch der Community. Bei den mock- und
fuzz-testing Tools, wird dieser Aspekt aus gelassen, da es sich bei diesen Tools
vorwiegend um Erweiterungen zu den unit-testing Tools handelt.
\newline

Zum Vergleich der einzelnen unit-testing Tools untereinander werden diese auf
den in Listing \ref{listing:base:my_module} abgebildeten Code angewendet. Da
sich diese Arbeit auch mit \gls{mock}-testing Tools beschäftigt, wurde auch für
diese ein Modul geschrieben, welches zu testen ist.
Der Code dazu befindet sich in Listing \ref{listing:base:my_mock_module}. Durch
die Anwendung des jeweiligen Tools auf diesen Code, ist es möglich die Aspekte
dieser untereinander zu vergleichen.

Das Modul aus Listing \ref{listing:base:my_module} für die unit-testing Tools
enthält eine selbst geschriebene Implementierung der Funktion \lstinline|pow()|,
welche eine Zahl \lstinline|a| mit einer Zahl \lstinline|b| quadriert. Um die
Komplexität zu erhöhen wurde eine in-memory Datenbank implementiert, welche
eine Items Tabelle enthält. Jedes Item hat eine ID (\lstinline|id|), einen Namen
(\lstinline|name|), einen Lagerplatz (\lstinline|storage_location|) und eine
Anzahl der vorrätigen Items \lstinline|amount|. Dazu besitzt jedes Item eine
Methode \lstinline|do_something()| welche eine externe Funktion/Methode,
die jedoch noch nicht existiert \lstinline|do_something_which_does_not_exist()|,
aufruft.

Für die \gls{mock}-testing Tools wurde eine Klasse geschrieben, welche 
Funktionen besitzt, durch ihren Namen ausdrücken, welche Funktionalität sie 
eigentlich haben sollten. Um jedoch testen zu können, ob das jeweilige Tool 
diese Methode \gls{stub}en kann, gibt nicht jede dieser Methoden das gewünschte 
Ergebnis zurück. Das Tool soll dafür sorgen, dass das gewünschte Ergebnis 
erfolgt. Selbstverständlich ist realer Code viel komplexer als in diesem Listing
(\ref{listing:base:my_mock_module}) dargestellt, jedoch reicht dieser Code aus,
um viele Tools an die Grenzen Ihrer Funktionalität zu bringen. So muss in
\lstinline{call_internal_functin_n_times} überprüft werden, ob sie \lstinline{n}
mal aufgerufen wurde. Auch \lstinline{call_helper_help} wird nicht ohne 
weiteres funktionieren, da die Methode der Klasse \lstinline{Helper} noch nicht 
implementiert wurde. Auch die Methode \lstinline{return_false_filepath}
bringt das ein oder andere Tool an die grenzen seiner Funktionalität, da eine 
externe Methode aus der STDLIB ersetzt werden muss, wobei dies nicht global 
geschehen darf.

Für die Fuzz-testing Tools werden keine Beispiele geschrieben, da diese meist 
auf komplexen Beispielen basieren um Sinn zu ergeben. Stattdessen werden im 
jeweiligen Kapitel kleinere Beispiele genannt, wie das Tool ein zu setzen ist 
oder es wird auf Beispiele aus der jeweiligen Dokumentation verwiesen.

Verfügt ein Tool über keinen Test-runner, so wird der von der STDLIB gestellte 
Runner \lstinline{unittest} verwendet.
    % !TeX root = ../bachelor.tex
\section{Python Test-Tool Analyse}\label{python-tools}
In diesem Kapitel werden die einzelnen Tools anhand der in Kapitel 
\ref{methodik} beschriebenen Methodik analysiert und verglichen. Am Ende des 
Kapitel wird für jede Kategorie (unit-, \gls{mock}- und \gls{fuzz}) eine 
Empfehlung anhand der jeweiligen Analyse gegeben, welches Tool die beste 
Eignung für die Anwendung von TDD hat.

% !TeX root = ../bachelor.tex
\subsection{Tools der Standard Bibliothek}\label{python-tools:stdlib}

Die Standard Bibliothek von Python bietet zwei verschiedene Test-Tools (\cite{wiki.python:PythonTestingToolsTaxonomy}).
Zum einen
\href{http://pyunit.sourceforge.net/pyunit.html}{unittest}\footnote{http://pyunit.sourceforge.net/pyunit.html}
und zum anderen
\href{https://docs.python.org/3/library/doctest.html}{doctest}\footnote{https://docs.python.org/3/library/doctest.html}.
Diese beiden Tools sind in ihrem Umfang bereits so vielseitig, dass es für
einen Entwickler unproblematisch ist, eine hohe Testabdeckung eines Programms
oder einer Bibliothek zu erreichen.
\newline
\\
Beide Tools zählen zu den "`Unit Testing Tools"' 
(\cite{wiki.python:PythonTestingToolsTaxonomy}),
auf Deutsch Modulare Testwerkzeuge, mit deren Hilfe die einzelnen Module eines 
Programms getestet werden können. In einem Programm oder einer Bibliothek wären 
dies die einzelnen Funktionen und Methoden.

% !TeX root = ../bachelor.tex
\subsubsection{unittest}\label{python-tools:unittest}

Unittest

% !TeX root = ../bachelor.tex
\subsubsection{doctest}\label{python-tools:doctest}

"`Das doctest Modul sucht nach Textstücken, die wie interaktive Python-Sitzungen
aussehen, und führt diese Sitzungen dann aus, um sicherzustellen, dass sie genau
wie gezeigt funktionieren."'\footnote{Übersetzt aus dem Englischen}
(\cite{docs.python:doctest}). Diese Textstücke müssen sich in Kommentaren
befinden, da sie sonst von Python als Code interpretiert werden.

\lstinline{doctest} bietet dem Nutzer eine Möglichkeit mithilfe von einem Test
den Code zu dokumentieren. Da \lstinline{doctest} lediglich Code ausführt wie in
einer interaktiven Python-Sitzung, ist auch nur das möglich aus zu führen was im
Code geschrieben ist. Test Fälle oder gar \Glspl{mock} sind hierbei nur bedingt
realisierbar, \Glspl{fixture} hingegen sind teilweise realisierbar in Form von
Code der vor dem Test ausgeführt wird.

\lstinline{doctest} selbst nutzt keine Assert Funktionen oder Methoden,
stattdessen wird der Output eines Befehls überprüft. So ergibt eine Funktion
\lstinline{return_3(s=None, i=None)} bei \lstinline{s=True} eine \lstinline{'3'}
und bei \lstinline{i=True} eine \lstinline{3}. Dieses Beispiel ist in Listing
\ref{listing:doctest:example} zu sehen, um einen Output zu erzeugen wurde noch
ein fehlerhafter Test hinzu gefügt. Dieser Output ist in Listing
\ref{listing:doctest:example_output} zu sehen.

Möchte der Entwickler einen etwas größeren Test schreiben, so kann er sich einer
Funktionalität von \lstinline{doctest} bedienen, die es Ihm ermöglicht Tests in
eine externe Datei zu verlagern. Dabei wird die Datei als ein \Gls{docstring}
behandelt, wodurch sie keine \lstinline{"""} benötigt. Um jedoch Code in diese
Datei ausführen zu können muss das zu testende Modul importiert werden.

Durch das schreiben der Tests in den \Glspl{docstring} werden die Module in
denen Tests geschrieben wurden schnell sehr unübersichtlich und lang, wenn der
Entwickler große Tests schreibt. Zwar wird durch externe Textdateien Abhilfe
geschaffen, dennoch sind zu lange und komplexe doctests schwer zu lesen und
nach zu vollziehen.
\newline

Das Listing \ref{listing:doctest:advanced} zeigt den Code aus
\ref{listing:base:my_module} mit \Glspl{docstring} versehen. Dieser Test wurde
mithilfe der \lstinline{main} Funktion von Python ausgeführt. Der in Listing
\ref{listing:doctest:advanced_text} ausgelagerte Test, wurde mithilfe des CLIs
von \lstinline{doctest} ausgeführt. Da beide Test keine Fehler werfen existiert
auch kein Output für diese Tests.
Die einzige Möglichkeit hier einen Output zu bekommen wäre mit \lstinline{-v} im
CLI oder \lstinline{verbose=True} im Funktionsaufruf. Die dort dargestellten
Informationen sind lediglich welche Kommandos ausgeführt wurden, was erwartet
wurde und das, dass erwartete eingetroffen ist. Ein kleines Beispiel ist in
Listing \ref{listing:doctest:output} zu sehen.

Da \lstinline{doctest} selbst keine Test Fälle unterstützt besitzt
\nameref{python-tools:unittest} eine Integration für \lstinline{doctest},
diese ermöglicht es Tests aus Kommentaren sowie Textdateien in unittest zu
integrieren und zu gliedern. Wie \nameref{python-tools:unittest}, ist auch
doctest in der STDLIB, wodurch keine externen Abhängigkeiten geladen werden
müssen.
\newline

Im Bezug auf die Anwendbarkeit von \lstinline{unittest} lässt sich sagen, dass
nicht alles zur Verfügung steht um TDD zu betreiben. Tests können nicht vor der
eigentlichen Funktionalität geschrieben werden und  \Glspl{fixture} und
\Glspl{mock} werden nicht unterstützt.

Die Anforderungen für \lstinline{doctest} sind sehr gering. Das testen geschieht
mithilfe der interaktiven Python-Sitzung, welche jedem Python Entwickler bekannt
ist. Die Tests selbst sind aufrufe der geschriebenen Funktionen und Methoden und
ein Abgleich des Outputs mit dem Erwarteten Wert. Demnach kann ein Entwickler
sehr Effizient mit \lstinline{doctest} umgehen ohne viel wissen zu benötigen.

\lstinline{doctest} selbst bietet keinerlei Komplexität von sich aus,
stattdessen ist die Komplexität die möglich ist vom Entwickler abhängig, da
dieser die Tests komplett selber schreiben muss. Je nach zu testender
Funktion/Methode kann dies abhängig des Test Aufwands schnell unübersichtlich
werden, wenn viel Code abseits des eigentlichen Tests benötigt wird.

Als Erweiterung bestehen lediglich die Integration in
\lstinline{unittest}\footcite{docs.python:doctest} sowie
\lstinline{pytest}\footcite{docs.pytest.org:4.4}.
\newline
\newline
\\
\textbf{Muss nach Zusammenfassung}
Da sich mit Doctest leider Funktionen nicht präparieren lassen, ist dieses
Test Modul für die alleinige Anwendung in TDD nicht nutzbar. Das Tool bietet
keine Möglichkeiten Objekte zu \gls{mock}en wodurch Tests an nicht
implementieren Methoden und Funktionen scheitern. Auch \Glspl{fixture} sind nur
bedingt realisierbar und benötigen viel Code der dupliziert werden muss, da
dieser vor und nach jedem Test geschrieben werden muss.

% !TeX root = ../bachelor.tex
\subsection{Tools abseits der Standard Bibliothek}\label{python-tools:extlib}
Abseits der Standard Bibliothek gibt es einige Tools, deren Nutzung von
Vorteil gegenüber der STDLIB ist. Diese werden unter diesem Punkt aufgeführt.
Auf \url{https://wiki.python.org/moin/PythonTestingToolsTaxonomy} werden viele
externe Tools gelistet. Jedoch scheinen viele inaktiv zu sein, da ihre
\glspl{commit} teilweise mehr als ein Jahr zurück liegen. Dies ist weder für
eine Bibliothek noch für ein Tool ein gutes Zeichen, da sich die Anforderungen
stetig ändern und niemals alle Fehler behoben und Features implementiert sein
können.

Manche der dort aufgelisteten Tools sind bereits oder werden gerade in andere
Tools integriert. Dies ist nötig, da ein Tool eine Erweiterung für ein anderes
war und die Entwickler die Änderungen angenommen haben oder um die Tools zu
verbessern und mehr Entwickler zur Verfügung zu haben. Selbstverständlich
können auch andere Gründe dafür verantwortlich sein.

Da sowohl Software als auch Programmiersprachen sich stetig weiter entwickeln,
werden in dieser Arbeit nur jene Tools behandelt, die sich diesen Entwicklungen
anpassen. Diese Anpassung zeigt sich durch die Unterstützung der aktuellsten
Version von Python oder durch neue Innovationen, sowie \Gls{bug}-fixes. Aus
diesem Grund werden Tools, deren letzter \Gls{commit} weiter als ein Jahr
zurück liegt, hier nicht behandelt.

Lässt man zusätzlich die Erweiterungen von Tool zunächst außen vor, so bleibt
lediglich
\href{https://github.com/pytest-dev/pytest/}{pytest}\footnote{https://github.com/pytest-dev/pytest/}
übrig (\cite{wiki.python:PythonTestingToolsTaxonomy}).
Tools, wie
\href{https://www.reahl.org/docs/4.0/devtools/tofu.d.html}{reahl.tofu}\footlabel{reahl.tofu}{https://www.reahl.org/docs/4.0/devtools/tofu.d.html}
oder
\href{https://pypi.org/project/zope.testing/}{zope.testing}\footlabel{zope.testing}{https://pypi.org/project/zope.testing/},
sind zwar mehr oder weniger aktiv, da sie beide auf die neusten Python
Versionen patchen, jedoch bieten sie sonst keinen Mehrwert in ihren Updates.
\href{https://www.reahl.org/docs/4.0/devtools/tofu.d.html}{reahl.tofu}\footref{reahl.tofu}
selbst ist auch nur eine Erweiterung für bestehende Test Tools, wie zum
Beispiel lstinline{pytest}, weshalb es hier nicht behandelt wird. So wird
\href{https://pypi.org/project/zope.testing/}{zope.testing}\footref{zope.testing}
auch nicht behandelt, da dieses Tool wie bereits beschrieben, nicht aktiv weiter
entwickelt wird, da es eher,
\href{http://www.zope.org/en/latest/}{zope}\footnote{http://www.zope.org/en/latest/}
Applikationen und nicht Python Applikationen testet.

Der Test-Runner
\href{https://pypi.org/project/nose/1.3.7/}{nose}\footnote{https://pypi.org/project/nose/1.3.7/},
der eine Erweiterung zu \nameref{python-tools:unittest} bietet, ist nicht mehr
in Entwicklung, dennoch haben sich ein paar Enthusiasten des Tools
zusammengeschlossen und
\href{https://pypi.org/project/nose2/}{nose2}\footlabel{nose2}{https://pypi.org/project/nose2/}
geschrieben. Da \href{https://pypi.org/project/nose2/}{nose2}\footref{nose2},
wie sein Vorgänger eine Erweiterung zu \nameref{python-tools:unittest} 
darstellt, wird es hier nicht analysiert.

Auch sehr bekannt ist
\href{https://pypi.org/project/testtools/}{testtools}\footnote{https://pypi.org/project/testtools/},
welches jedoch auch als Erweiterung zu \nameref{python-tools:unittest} gesehen 
werden muss. Hier findet ebenfalls keine weitere Entwicklung statt.
\newline
\\
Als weitere Kategorie werden \Gls{mock}-Tools geführt. Auch wenn fast alle
unittest-Tools integriertes \gls{mock}ing haben, lässt sich mit diesen
Tools meist mehr erreichen. Es wurden die gleichen Filterkriterien verwendet
wie bei den unit-testing Tools (\cite{wiki.python:PythonTestingToolsTaxonomy}).
Die demnach relevanten Tools sind:
\href{https://www.reahl.org/docs/4.0/devtools/stubble.d.html}{stubble}\footnote{https://www.reahl.org/docs/4.0/devtools/stubble.d.html},
\href{https://github.com/timbertson/mocktest/tree/master}{mocktest}\footnote{https://github.com/timbertson/mocktest/tree/master},
\href{https://github.com/bkabrda/flexmock}{flexmock}\footnote{https://github.com/bkabrda/flexmock},
\href{https://bitbucket.org/DavidVilla/python-doublex}{python-doublex}\footnote{https://bitbucket.org/DavidVilla/python-doublex}
und 
\href{https://github.com/ionelmc/python-aspectlib}{python-aspectlib}\footnote{https://github.com/ionelmc/python-aspectlib}.
\newline
\\
Als letzte Kategorie werden hier die \Gls{fuzz} Tools behandelt, da diese eine
gute Möglichkeit bieten, Code ausgiebig zu testen. Das wohl umfangreichste und
nach den obigen Kriterien einzige Tool ist
\href{https://github.com/HypothesisWorks/hypothesis}{hypothesis}\footnote{https://github.com/HypothesisWorks/hypothesis}
(\cite{wiki.python:PythonTestingToolsTaxonomy}).

% !TeX root = ../bachelor.tex
\subsubsection{pytest}\label{python-tools:pytest}\mbox{}
\newline
Das am 04. August 2009 in Version
1.0.0\footnote{https://github.com/pytest-dev/pytest/releases/tag/1.0.0}
veröffentlichte Tool \lstinline{pytest} (auch py.test genannt)
ist ein sehr umfangreiches und weit entwickeltes Tool. Seit 2009 wird das Tool
stets weiter entwickelt und vorangetrieben, wodurch es eine Menge an Features
gewonnen hat.
\noindent
Die Basis Features von \lstinline{pytest} sind folgende(\cite{docs.pytest.org:4.4}):
\begin{itemize}
    \item Simple \lstinline{assert} statements.
    \subitem Kein \lstinline{self.assert}
    \subitem Error Überprüfung mit \Gls{context}
    \item Informativer Output in Farbe
    \subitem Der gesamte Output kann angepasst werden
    \item Feature reiche \Glspl{fixture}
    \subitem Vordefinierte \Glspl{fixture} von \lstinline{pytest}
    \subitem Geteilte \Glspl{fixture} unter Tests
    \subitem Globale \Glspl{fixture} zwischen Modulen
    \subitem Parametrisierung von Test als \Glspl{fixture}
    \item Überprüfung von \lstinline{stdout} und \lstinline{stderr}
    \item Gliedern von Test in Fällen
    \subitem ... durch Markierung
    \subitem ... durch Nodes (Auswahl der Modul Abhängigkeit)
    \subitem ... durch String Abgleich der Funktions-Namen
\end{itemize}
\noindent
Wie anhand der Basis Features erkennbar ist, bietet \lstinline{pytest} einiges
um Tests zu schreiben und aus zu führen. So ist mit den gebotenen
\Glspl{fixture} bereits eine Voraussetzung für TDD erfüllt, da
\lstinline{pytest} viele verschiedene Arten bietet diese zu benutzen.
\Glspl{fixture} werden in \lstinline{pytest} allerdings nicht mit
\lstinline{setUp} und \lstinline{tearDown} geschrieben, sondern werden mithilfe
eines \glspl{decorator} markiert und den Test Funktionen als Parameter
übergeben. Um die \lstinline{setUp} und \lstinline{tearDown} Funktionalität zu
bekommen muss lediglich das keyword \lstinline{yield} verwendet werden, die
\Gls{fixture} wird dann den Code bis zum \lstinline{yield} ausführen und nach
der Funktion den Rest nach \lstinline{yield} ausführen.
\newline
\\
Diese lassen sich durch zusätzliche Parameter weiter anpassen. Diese
Funktionen sind in Kapitel fünf nach \cite{docs.pytest.org:4.4} beschrieben.

Des weiteren bietet \lstinline{pytest} auch \gls{mock}ing, so lässt sich beispielsweise mit
\lstinline{monkeypatch.setattr()} der return Wert einer Funktion ersetzen. Dies wird
in \lstinline{pytest} automatisch am ende der Funktion rückgängig gemacht, wodurch der Entwickler
sich mehr auf das eigentliche Testen konzentrieren kann.
\newline
\\
Auch hier bietet \lstinline{pytest} weitere Möglichkeiten \Glspl{mock} zu verwenden, dazu hat
\cite{docs.pytest.org:4.4} in Kapitel sieben einige Worte geschrieben.

Ein weiteres sehr praktisches Feature von \lstinline{pytest} ist die Gliederung in Fälle. Dies ist bei \lstinline{pytest}
sehr fein einstellbar, so lässt sich wie in den Basis Features bereits beschrieben ein Test mit
einer oder mehreren Markierungen versehen wodurch eine Gliederung nach Markierung entsteht. Auch
durch das selektieren bestimmten Wörtern lassen sich Tests nach ihrem Namen gliedern, so entsteht
einerseits eine Gliederung zur Ausführung der Tests und andererseits eine Gliederung für die
Entwickler die sie selbst im Code sehen können. Als letzte alternative lassen sich Tests anhand
ihrer Module ausführen, dies geschieht durch die Angabe der Module. So würde
\lstinline{test_datei.py::TestKlasse::test_methode} den Test \lstinline{test_methode} der Klasse
\lstinline{TestKlasse} in der Datei \lstinline{test_datei.py} ausführen.
\newline
\\
Diese Features werden in Kapitel sechs von \cite{docs.pytest.org:4.4} beschrieben.

Selbstverständlich bietet \lstinline{pytest} noch weitere Features jedoch sind diese nicht zwangsläufig notwendig
um TDD zu betreiben und sind mehr ein nice to have Feature als wirklich benötigt. Für eine
Vollständige Auflistung und Erklärung aller Features kann jederzeit unter
\url{https://docs.pytest.org/en/latest/} die aktuellste Version der Dokumentation abgefragt werden.

Da \lstinline{pytest} nicht in der STDLIB ist, muss es mit einem \Gls{paket} installiert werden.
Dabei werden für \lstinline{pytest 4.4.0} die in Listing \ref{listing:pytest:requirements}
gezeigten Abhängigkeiten installiert. Demnach benötigt \lstinline{pytest} sechs externe Abhängigkeiten
zusätzlich zu sich selbst.
Demnach bietet \lstinline{pytest} alles um TDD anwenden zu können und die sieben zusätzlichen Abhängigkeiten
die ein Entwickler bei der Nutzung von \lstinline{pytest} eingeht sollten keinen Entwickler davon abhalten
\lstinline{pytest} und seine Features zu genießen.
\newline
\\
Da \lstinline{pytest}, wie bereits erwähnt keine neuen assert Methoden hinzufügt lässt sich sehr schnell und
einfach ein Test schreiben. Selbst die Nutzung von \Glspl{fixture} ist in \lstinline{pytest} sehr einfach, da
lediglich eine Funktion geschrieben werden muss die mit \lstinline{@pytest.fixture} markiert wurde
und der Test Funktion oder Methode als Parameter übergeben wird. Den Rest erledigt \lstinline{pytest} selbst
im Hintergrund. Genauso einfach gestaltet sich die Nutzung von \Glspl{mock}.

Die Effizienz, die bei der Nutzung von \lstinline{pytest} entsteht ist demnach sehr hoch. Denn der Entwickler
muss nicht wirklich etwas neues dazu lernen um Tests zu schreiben oder zu präparieren. Genauso
leicht kann ein Entwickler den Output von \lstinline{pytest} auswerten, da dieser erstens, in Farbe ist, was
die Lesbarkeit deutlich erhöht, zweitens sehr gut gegliedert ist und Wichtige Strukturen klar
darstellt und drittens nach den Wünschen der Entwickler sich gestalten und verbessern lässt.
\newline
\\
Durch die eben genannten Features im Bezug auf das schreiben von Tests mit assert, sowie
\Glspl{fixture} und \Glspl{mock} lässt sich sagen das \lstinline{pytest} sehr viel Funktionalität bietet
wobei es trotzdem Struktur im Code bietet und komplexe Features einfach ermöglicht. Demnach
kann ein Entwickler mit wenig Code viel Testing Funktionalität erstellen.
\newline
\\
Sollten dem Entwickler die Features von \lstinline{pytest} nicht ausreichen, so findet man unter
\url{http://plugincompat.herokuapp.com/} eine Liste von 618 (Stand: 2. April 2019) Erweiterungen
für \lstinline{pytest 4.3.0}. \lstinline{pytest} selbst kann allerdings auch als Erweiterung zu unittest
genutzt werden, indem ein Entwickler \lstinline{pytest} zum ausführen der unittests verwendet. Dadurch bietet
sich dem Entwickler ein verbesserter Output des Test Ergebnisses.
\newline
\\
Im Listing \ref{listing:pytest:advanced} und \ref{listing:pytest:advanced_output} befindet sich der
Code und der Output zum Test von dem in Listing \ref{listing:base:my_module} definierten Modul.
% !TeX root = ../bachelor.tex
\subsubsection{\Gls{mock}ing Tools}\label{python-tools:extlib:mock}

Um eine bessere Gliederung zur Verfügung zu stellen, werden die \Gls{mock}ing
Tools in diesem Kapitel gegliedert.

Wie in \nameref{python-tools} bereits beschrieben, werden hier Tools analysiert,
die eine Erweiterung für Test Tools im Bezug auf verbessertes \gls{mock}ing zur 
Verfügung stellen.

\input{bachelor/3_2_2_1_stubble}
\input{bachelor/3_2_2_2_mocktest}
\input{bachelor/3_2_2_3_flexmock}
\input{bachelor/3_2_2_4_python_doublex}
% !TeX root = ../bachelor.tex
\subsubsection{Fuzz-testing Tools}\label{python-tools:extlib:fuzz}

Zusätzlich zu den hier behandelten Test runnern und den \gls{mock}ing Tools 
werden Fuzz-testing Tools behandelt. Unter Fuzz-testing 
versteht man das Testen eines Programms oder einem Modul mit zufälligen Werten. 
Diese Werte können komplett zufällig sein, aber auch einem gewissen Format 
oder einem Typ entsprechen.

Fuzz-testing funktioniert dabei anders als die gewohnten Methoden des Testens. 
Normalerweise werden im Test Daten präpariert, die im Test ausgeführt und 
danach validiert werden, sodass das Ergebnis richtig ist. Mit Fuzz-testing 
sieht das ganze etwas anders aus: zuerst wird festgelegt, welchem Schema die 
Daten entsprechen müssen (Beispiel: Es wird eine IP erwartet, 
\lstinline|[1-9]{1,3}.[1-9]{1,3}.[1-9]{1,3}.[1-9]{1,3}|), danach wird mit den 
Daten der Test ausgeführt und schlussendlich validiert 
(\cite{hypothesis:doc:4.18.0}).

Der Unterschied ist, dass der Entwickler sich selbst keine Daten ausdenken 
muss. Im genannten Beispiel müsste der Entwickler sich verschiedene IP Adressen 
ausdenken, die entweder richtig oder falsch sind und danach diese überprüfen. 
Fuzz-testing übernimmt das Ausdenken der Daten und sorgt dafür, dass diese der 
Spezifikation entsprechen. So kann eine viel größere Anzahl an Tests mit 
erheblich weniger Aufwand auf einem Modul oder Programm ausgeführt werden.

Schlägt ein Test mit Daten, die per  Spezifikation richtig sind, fehl, so wird 
der Wert der Daten für später gespeichert und, sofern möglich, versucht mit 
ähnlichen Daten dieser Fehler zu erreichen. Dadurch hat ein Entwickler am Ende 
der Tests Daten, die fehlgeschlagen sind. Mit diesen kann er die weiteren Tests 
befüllen um so  zu verhindern, dass diese Daten zu Fehlern führen. Im 
Normalfall übernimmt dies allerdings das Tool, wodurch der Entwickler sich nur 
darum kümmern muss, dass die Tests erfolgreich verlaufen.

Von Zeit zu Zeit entsteht dadurch ein Katalog an Daten für verschiedene Tests, 
der vorgibt bei welchen Daten ein Test fehl geschlagen ist. Dadurch wird 
verhindert, dass bei der Änderung des Codes dieser Fehler wieder auftritt. 
Führt man  dies über den gesamten Entwicklungsprozess hinweg durch, so erhält 
man am Ende ein sehr ausgiebig getestetes Programm.

Diese Methode des Testens wird auch "`property based testing"', also 
Eigenschaftsbasierte Tests genannt. In dieser Arbeit, wird allerdings weiterhin 
der englische Begriff Fuzztesting genutzt.

\input{bachelor/3_2_3_1_hypothesis}

% !TeX root = ../bachelor.tex
\subsection{Vergleich der Tools}\label{python-tools:vergleich}

Die in Kapitel \ref{python-tools:stdlib} und \ref{python-tools:extlib}
analysierten Tools, werden anhand der Ergebnisse aus der Anwendung des
jeweiligen Tools in diesem Kapitel verglichen. Der Vergleich wird in unit- und
\gls{mock} Tools gegliedert. Die \gls{fuzz} Tools werden hier
nicht verglichen, da nur ein Tool analysiert wurde. Am Ende dieses Kapitels 
wird für die jeweilige Kategorie (unit- und mock-testing Tools) eine Empfehlung 
ausgesprochen, die anhand der Ergebnisse des Vergleichs festgelegt wird. Anhand 
dieser Empfehlungen werden in Kapitel \ref{python-tools:kombination} 
verschiedene Möglichkeiten aufgezeigt, die Tools untereinander zu kombinieren 
um mehr Funktionalität für eine mögliche Anwendung zu bekommen.

\subsubsection{Unit-testing Tools}\label{python-tools:vergleich:unit}
Die unit-testing Tools belaufen sich auf drei Tools, \lstinline{unittest},
\lstinline{doctest} und \lstinline{pytest}. Die ersten beiden Tools sind beide
in der STDLIB, während \lstinline{pytest} mit sechs weiteren Paketen betrieben
werden muss. Während \lstinline{unittest} und \lstinline{pytest} sich sehr
ähnlich sind, ist \lstinline{doctest} vollkommen anders im Bezug auf Anwendung
und gegebener Funktionalität. \lstinline{doctest} unterstützt keine Testfälle
oder Testumgebungen, stattdessen wird im \gls{docstring} der jeweiligen
Funktion oder Methode der Test geschrieben. Zwar ist es möglich diesen in eine
externe Datei aus zu lagern um so einen zentralen Punkt zu schaffen, an dem sich
alle Tests befinden (Testumgebung/ Testfälle), jedoch ist hier die Datei das
Konstrukt, dass die Gliederung ermöglicht. Sowohl \lstinline{unittest} als auch
\lstinline{pytest} bieten hier die Möglichkeiten, Tests in Testumgebungen
(Klassen) und Testfällen (Funktionen/ Methoden) zu gliedern. \lstinline{pytest}
bietet zusätzlich auch noch die Möglichkeit diese Tests zu markieren um eine
noch feinere Gliederung zu schaffen oder um zwischen Testumgebungen und 
Testfälle zu kombinieren. Im Fall der Gliederung und Verwaltung von Tests sind 
hier sowohl \lstinline{pytest} als auch \lstinline{unittest}, 
\lstinline{doctest} überlegen. Lässt man die \lstinline{unittest} Integration 
für \lstinline{doctest} außen vor, so ist \lstinline{unittest} in diesem Fall 
der Gewinner. Wird allerdings die Anwendung der \lstinline{unittest} 
Integration von \lstinline{doctest} in Betracht gezogen, so sind 
\lstinline{unittest} und \lstinline{doctest} auf einem Level, was Gliederung 
und Verwaltung von Tests betrifft.

Ein Feature, das eine Gliederung voraussetzt, sind die Fixtures. Da
\lstinline{doctest} keine Gliederung bietet, sind Fixtures für dieses Tool
nicht als Feature integriert. Es ist zwar möglich Fixtures manuell zu
schreiben, jedoch ist der Aufwand, der dafür betrieben werden muss, zu hoch, als
dass es sich rentieren würde. Die \lstinline{unittest} Integration ermöglicht es
zwar Fixtures zu nutzen, jedoch wird diese Integration außen vor
gelassen. \lstinline{unittest} und \lstinline{pytest} hingegen unterstützen
Testfixtures mit ihren \lstinline{setUp()} und \lstinline{tearDown()} Methoden.
Beide Tools bieten die Möglichkeit die Fixtures sowohl für jedem Testfall, jede
Testumgebung oder jedes Modul aus zu führen. \lstinline{pytest} bietet
zusätzlich noch weitere Funktionalität, sowie fertige Fixtures die dem
Entwickler bei der Nutzung zur Verfügung stehen. Diese wurden in Kapitel
\ref{python-tools:pytest} beschrieben und sind in der 
\href{https://docs.pytest.org/en/latest/fixture.html}{Dokumentation}\footnote{\url{https://docs.pytest.org/en/latest/fixture.html}}
zu finden. Im Bezug auf Testfixtures ist also \lstinline{pytest} das Tool
mit den meisten Features, gefolgt von \lstinline{unittest}. \lstinline{doctest}
scheidet bei dieser Kategorie vollkommen aus, da die Funktionalität gänzlich
fehlt.

Als nächstes integriertes Feature der unit-testing Tools sind \Glspl{mock} und
\Glspl{stub} zu betrachten. Um für TDD in Frage zu kommen sollte ein
unit-testing Tool \Glspl{stub} und \Glspl{mock} unterstützen. Eine ausführliche
Erklärung dazu ist in Kapitel \ref{python-tools} zu finden. Wie in Kapitel
\ref{python-tools:doctest} aus der Analyse von \lstinline{doctest} hervorgeht,
gibt es weder die Möglichkeit einen \Gls{mock} noch einen \Gls{stub} auf zu
setzen, weshalb \lstinline{doctest} in dieser Kategorie als unbrauchbar gesehen
werden kann. \lstinline{unittest} als Modul selbst, besitzt keine \gls{mock}ing
Funktionalität, jedoch ist es möglich, dass Submodul \lstinline{unittest.mock}
dafür zu verwenden. \lstinline{unittest.mock} ist zwar eine Erweiterung in Form 
eines Submoduls, weshalb es zu \lstinline{unittest} dazu gehört und in diesem 
Vergleich beachtet wird. Damit unterstützt \lstinline{unittest} sowohl das 
\Gls{mock}en, als auch das \Gls{stub}en. \lstinline{pytest} hingegen verfügt in 
seiner Basisausführung lediglich über die Möglichkeit des \Gls{stub}ens, nicht 
jedoch über die Möglichkeit ein Objekt zu \gls{mock}en. In diesem Falls ist 
also \lstinline{unittest} der klare Gewinner, gefolgt von \lstinline{pytest}.
\newline

Im Bezug auf die Effizienz der unit-testing Tools lässt sich folgendes
feststellen. \lstinline{doctest} ist zwar für jeden Entwickler einfach und
verständlich anwendbar, jedoch ist die Vorarbeit um Tests auf zu setzen im
Vergleich zu \lstinline{unittest} und \lstinline{pytest} höher. Sowohl
\lstinline{unittest} als auch \lstinline{pytest} benötigen das gleiche Maß an
Vorarbeit um Tests auf zu setzen, wobei der Aufwand das jeweilige Tool ein zu
setzen bei \lstinline{unittest} etwas höher ist, da es so viele verschiedene
\lstinline{assert} Methoden zu merken gibt. Die Auswertung der Tests erfolgt
bei allen drei unit-testing Tools auf dem Terminal. \lstinline{doctest} und
\lstinline{unittest} haben beide einen Output der voneinander fast nicht zu
trennen ist. Beide Tools haben den Nachteil gegenüber \lstinline{pytest}, dass
ihr Output nicht farbig, sondern monochrom ist.  Aus diesem Grund ist 
\lstinline{pytest}, \lstinline{unittest} und \lstinline{doctest} im Bezug auf
effiziente Anwendung vor zu ziehen, wobei \lstinline{unittest} noch etwas 
besser als \lstinline{doctest} ist.
\newline

Der Aspekt der Komplexität ist in zwei Unterkategorien gegliedert. So wird zum 
Einen die Fülle an Funktionen, die geboten werden und Komplexität liefern, und
andererseits die Menge an Code, die zu schreiben ist und die Anwendung und
Wartung Komplexer gestaltet, betrachtet. Im Falle der ersten Kategorie, erweist 
sich \lstinline{doctest} als nicht all zu komplex. Fast die Gesamte gebotene
Funktionalität bezieht sich auf die Ausführung, jedoch nicht das Schreiben der
Tests, welches alleine durch die Fähigkeit Python zu schreiben bestimmt wird.
\lstinline{unittest} und \lstinline{pytest} hingegen bieten dem Entwickler
einiges an Funktionalität um Tests zu schreiben und aus zu führen. Auch hier
kann \lstinline{unittest} nur mithalten, wenn das Submodul
\lstinline{unittest.mock} mit in die Wertung aufgenommen wird, da dieses
weitere Funktionalität bietet. \lstinline{unittest} und \lstinline{pytest}
haben viele gemeinsame Features, wovon lediglich das \gls{mock}ing in
\lstinline{pytest} nicht ohne Erweiterungen verfügbar ist. Ansonsten bietet
\lstinline{pytest} zusätzliche Funktionalität um Tests feiner ein zu
stellen oder um Randbedingungen ab zu decken. In dieser Kategorie ist
\lstinline{pytest} also etwas besser als \lstinline{unittest}. Beide Tools
sind \lstinline{doctest} bei Weitem im Bezug auf die gebotenen Funktionalitäten 
überlegen.

In der zweiten Kategorie, der Komplexität im Bezug auf Codelesbarkeit und
-umfang, erweisen sich \lstinline{unittest} und \lstinline{pytest} als 
gleichauf. Beide Tools haben einen ähnlichen Aufbau, mit den der Code 
geschrieben wird (vgl. Listing \ref{listing:unittest:advanced} und
\ref{listing:pytest:advanced}). Die Lesbarkeit ist bei beiden Tools gleich gut
und die Menge an Code, die benötigt wird, gleicht sich auch. Hier ist es dem
Entwickler überlassen zu entscheiden, welches Tool den schöneren Code schreibt.
\lstinline{doctest} hingegen produziert viel Code, der mit zunehmender
Funktionalität an Lesbarkeit verliert.Im Bezug auf die Komplexität der Tools
ist \lstinline{pytest} also knapp vor \lstinline{unittest}, gefolgt von
\lstinline{doctest}.
\newline

Als letztes ist die Erweiterbarkeit der unit-testing Tools zu vergleichen.
Beginnend mit \lstinline{doctest}, lässt sich sagen, dass es zwei sehr gute
Erweiterungen gibt, die das Tool erheblich verbessern. Diese Erweiterungen sind
die Integration in \lstinline{unittest} und \lstinline{pytest}. Erstere wird
von der Python Sprache selbst geboten und ist in \lstinline{unittest} verfügbar,
zweite wird von \lstinline{pytest} geboten und ist auch im Basis Paket
enthalten. Abgesehen von diesen zwei gibt es noch weitere Integrationen in
andere Toolsets, jedoch wurden diese hier wie in Kapitel \ref{methodik} 
beschrieben, nicht behandelt. Das Modul \lstinline{unittest} verfügt über 
einige Erweiterungen welche neue Funktionalität hinzu fügen oder alte 
Funktionalität verbessern. Das Tool mit den meisten Erweiterungen ist 
\lstinline{pytest}, welches derzeit 683 (Stand: 19. Juli 
2019\footnote{\url{http://plugincompat.herokuapp.com/}}) Erweiterungen
bietet, wovon die meisten sehr nützliche Funktionalität bieten.
\newline

Im Vergleich der unit-testing Tools ergibt sich also ein klarer Sieger:
\lstinline{pytest}. Das Tool ist in fast allen Analysierten Aspekten die beste
Wahl für den Einsatz mit TDD. Jedoch bietet \lstinline{unittest} für
Entwickler, die keine externen Abhängigkeiten möchten eine Alternative, die
durchaus mit \lstinline{pytest} mithalten kann.

\subsubsection{Mock-testing Tools}\label{python-tools:vergleich:mock}
Die analysierten \gls{mock}-testing Tools, sind \lstinline{stubble},
\lstinline{mocktest}, \lstinline{flexmock} und \lstinline{doublex}. Alle vier
bieten verschiedene Möglichkeiten \Glspl{mock} und/oder \Glspl{stub} ein zu
setzen. Da es sich hier um Tools handelt, deren Funktionalität sich um das
\Gls{mock}en von Objekten dreht, wird der erste analysierte Aspekt, die
Anwendbarkeit, auf die gebotene Auswahl an Funktionalität im Bezug auf
\Gls{mock}en und \Gls{stub}en betrachtet.

Das einzige Tool, das bei der Anwendung auf den zu testenden Code keine
Probleme mit sich brachte, ist \lstinline{mocktest}. Das Tool bietet im Aspekt 
der Anwendbarkeit sämtliche Features um effektiv \Glspl{mock} und \Glspl{stub} 
ein zu setzen. \lstinline{mocktest} bietet seinem Anwender eine deskriptive 
Schreibweise, mit der es problemlos möglich ist Tests zu schreiben. Des 
Weiteren wird eine Testisolation geboten, mit deren Hilfe es möglich ist, Tests 
besser voneinander ab zu schotten. Zusätzlich benötigt das Tool keine weiteren
Abhängigkeiten zu sich selbst. Direkt danach wäre \lstinline{stubble} das  
nächstbeste Tool im Aspekt der Anwendbarkeit. So bietet das Tool zwar alles 
um  es für TDD zu verwenden, jedoch ist das \Gls{stub}en einer nicht 
existierenden Methode nur unter Einschränkungen möglich. Die Hauptfeatures von 
\lstinline{stubble} sind dabei die vorgefertigten \Glspl{stub}, welche 
nützliche Features für den Entwickler bieten. Um einen \Gls{stub} zu erstellen 
muss die zu \gls{stub}ende Klasse mit \lstinline{stubble} neu geschrieben 
werden. Zusätzlich benötigt \lstinline{stubble} eine weitere Abhängigkeit zu 
sich selbst. An dritter Stelle ist das Tool \lstinline{flexmock} welches,
abgesehen vom ersetzen einer nicht geschriebenen Methode um TDD an zu 
wenden, vieles bietet (vgl. Listing \ref{listing:flexmock:exception}). Wäre 
dieses Defizit nicht, würde das Tool im Aspekt der Anwendbarkeit sehr weit oben 
stehen. Flexmock vereint die \Glspl{mock} und \Glspl{stub} in einer Klasse, 
wodurch die gesamte Funktionalität mit einem Objekt geboten wird. Die Schreib 
weise der Tests ist hierbei, wie bei \lstinline{mocktest}, deskriptiv. Auch die 
benötigten Abhängigkeiten von \lstinline{flexmock} sprechen für das Tool, da 
keine weiteren zu sich selbst benötigt werden. Als letztes Tool im Aspekt der
Anwendbarkeit ist \lstinline{doublex} zu betrachten. Das Tool bietet zwar
einiges an Funktionalität um zu \gls{mock}en oder zu \gls{stub}en, jedoch ist
die Anwendung dieser Features für TDD nicht nutzbar. Es werden vier
verschiedene Interfaces geboten, deren Funktionalität sich überschneidet und
die jeweils ein Objekt hervorbringen, das ein Duplikat vom Zielobjekt ist.
Die daraus resultierenden Anwendungszwecke bieten wenig Nutzbarkeit für TDD.
Hinzu kommt, dass \lstinline{doublex} drei weitere Abhängigkeiten benötigt und
anscheinend nicht mehr aktiv weiter entwickelt wird.

Demnach ist im Aspekt der Anwendbarkeit \lstinline{mocktest} also
\lstinline{stubble} und \lstinline{flexmock} vor zu ziehen. \lstinline{doublex}
fällt aufgrund dessen, dass es nicht für TDD geeignet ist und derzeit nicht
aktiv weiter entwickelt wird, komplett raus. Aus diesem Grund wird dieses Tool
im weiteren Verlauf des Vergleichs nicht mehr berücksichtigt.
\newline

Als nächstes ist die Effizienz der jeweiligen Tools zu betrachten. Dabei spielt
vor allem der benötigte Aufwand das Tool ein zu setzen eine Rolle. Nebenbei
wird dabei auch noch auf die benötigte Vorarbeit geachtet, die geleistet werden
muss um das Tool an zu wenden. Sowohl bei \lstinline{mocktest} als auch bei
\lstinline{flexmock} ist der zu betreibende Aufwand sehr gering. Beide Tools
haben ein ähnliches Interface mit dem deskriptiv beschrieben wird, was ein
Objekt zu machen hat. Der dabei entstehende Aufwand ist sehr gering, da sich
der Code so schreibt, wie ein Entwickler sich die Funktionalität vorstellt.
Beide Tools benötigen die gleiche Menge an Vorarbeit und sind demnach im Bezug
auf die Effizienz mit der sie angewendet werden können gleich auf. Lediglich
\lstinline{stubble} ist in diesem Aspekt etwas schlechter als die beiden
anderen Tools. So ist die \gls{mock}ing Funktionalität ausschließlich mit einem
vorgefertigten \Gls{stub} möglich, wodurch dieser erst verstanden werden muss,
bevor er effektiv eingesetzt werden kann. Zusätzlich muss jedes zu
\gls{stub}ende Objekt überschrieben werden, wodurch die benötigte Vorarbeit 
sehr hoch ausfällt.

Demnach sollte im Aspekt der effizienten Anwendung der \gls{mock}-testing Tools 
\lstinline{stubble} nicht verwendet werden. Es ist dem Entwickler überlassen ob 
er in diesem Aspekt das Interface von \lstinline{mocktest} oder 
\lstinline{flexmock} lieber hat.
\newline

Zuletzt ist die Komplexität der \gls{mock}-testing Tools zu betrachten. Dabei
wird zum einen auf die gestellten Funktionalitäten abseits der Basis
Funktionalität geachtet und zum anderen auf die Lesbarkeit des entstehenden
Codes des Tools. Auch hier sind \lstinline{mocktest} und \lstinline{flexmock}
auf einem Level, was die zusätzliche Funktionalität betrifft. Jedoch hat
\lstinline{mocktest} etwas mehr zusätzliche Funktionalität gegenüber
\lstinline{flexmock} geboten (vgl. Kapitel \ref{python-tools:mocktest} und
\ref{python-tools:flexmock}). \lstinline{stubble} bietet auch in diesem Aspekt
weniger als die beiden Tools. Abseits der Basisfunktionalität wird von
\lstinline{stubble} nichts weiter geboten. Im Bezug der Lesbarkeit sieht es
ähnlich aus. \lstinline{mocktest} und \lstinline{flexmock} befinden sich auf
einem Level, während \lstinline{stubble} durch viel Code an Lesbarkeit
verliert. Im Fall \lstinline{mocktest} gegen \lstinline{flexmock} ist es wieder
dem Anwender überlassen, welches der beiden Tools lesbarer ist als das andere.

Im Bezug auf Komplexität geben sich \lstinline{mocktest} und
\lstinline{flexmock} fast nichts, jedoch hat \lstinline{mocktest} durch die
zusätzliche Funktionalität hier die Nase vorne. \lstinline{stubble} bleibt auch
hier auf dem letzten Platz.
\newline

Im gesamten Vergleich ist \lstinline{mocktest} durch seine Features und sein
Interface an erster Stelle, wenn es um den Einsatz von \gls{mock}-testing Tools
geht. Der zweite Platz ist schwer zu wählen, da \lstinline{flexmock} mit seiner
fehlenden Funktionalität nicht implementierte Methoden hinter
\lstinline{stubble} liegt, das jedoch auch nur bedingt diese Anforderung
erfüllt. In allen anderen Aspekten ist \lstinline{flexmock} \lstinline{stubble}
voraus. Aus diesem Grund erhält \lstinline{flexmock} den zweiten Platz der
\gls{mock}-testing Tools und \lstinline{stubble} den dritten. Auf dem vierten
wäre \lstinline{doublex}, jedoch ist dieses Tool für die Anwendung von TDD
ungeeignet, weshalb es keine Platzierung erhält.
% !TeX root = ../bachelor.tex
\subsection{Kombinierung von Tools}\label{python-tools:kombination}
Um eine verbesserte Funktionalität zu bekommen, ist es ratsam verschiedene
Tools miteinander zu kombinieren. Die in diesem Kapitel aufgezeigten
Möglichkeiten zur Kombinierung ergeben sich aus den analysierten Features der
jeweiligen Tools. Dabei wird zuerst geprüft ob die beiden, am besten geeigneten
Tools aus unit- und \gls{mock}-testing Tools miteinander verwendet werden
können, um so einen Mehrwert zu erhalten.

Bei den unit-testing Tools ist \lstinline{doctest} das am wenigsten geeignete
Tool für die alleinige Anwendung mit TDD. Dennoch bietet das Tool Features, die
für die Anwendung von TDD nützlich sind. So kann \lstinline{doctest} als
testende Dokumentation gesehen werden. Manche Funktionen oder Methoden sind in
ihrem Funktionsumfang so gering, dass ein ganzer Test als zu viel erscheint.
Bei diesen Funktionen oder Methoden kann beispielsweise ein
\lstinline{doctest} geschrieben werden, der diese kleine Funktionalität prüft.
Des-weiteren kann \lstinline{doctest} als ergänzendes Tool für die
Dokumentation des Codes gesehen werden, denn die Dokumentation ist ein
wichtiger Bestandteil eines Projekts (\cite{python.org:PEP8}).
\lstinline{doctests} können dafür sorgen, dass die Dokumentation besser
verstanden wird und eindeutiger ist. Aus diesen Gründen kann
\lstinline{doctest} immer als ergänzendes Tool in jede Entwicklung aufgenommen
werden.

Da das Tool \lstinline{hypothesis} alleiniger Vertreter seiner Kategorie, der 
\gls{fuzz} Tools ist, ist es das empfohlene Tool, wenn \gls{fuzz} benötigt wird.
Da \lstinline{hypothesis} mit jedem anderen Tool verwendet werden kann,
empfiehlt es sich, dieses so oft wie es möglich ist ein zu setzen um eine
bessere Testabdeckung zu erreichen.

Demnach ist es ratsam sowohl \lstinline{doctest} als auch 
\lstinline{hypothesis} immer ein zu setzen wo es möglich ist.
\newline

Grundsätzlich lässt sich ein jedes unit-testing Tool mit jedem  
\gls{mock}-testing Tool kombinieren. Der dadurch erhaltene Mehrwert ist meist 
beachtlich gegenüber der Verwendung der Tools alleine. Es ist jedem Entwickler 
überlassen, welches Tool er mit welchem kombiniert, dennoch bietet es sich an, 
\lstinline{pytest} zusammen mit \lstinline{mocktest} zu verwenden. Wenn dazu 
noch \lstinline{doctest} und \lstinline{hypothesis} benutzt wird, erhält der 
Entwickler ein Toolset, das ihm alles an Funktionalität bietet, was er 
benötigt. Ist es erwünscht \lstinline{unittest} statt \lstinline{pytest} 
zu verwenden, so kann dies auch getan werden. Jedoch zeigt die Analyse, dass 
der Einsatz von \lstinline{pytest} dem Einsatz von \lstinline{unittest} vor zu 
ziehen ist.
    % !TeX root = ../bachelor.tex
\section{Diskussion}\label{diskussion}
Um für die Anwendung von TDD die passenden Tools zu finden, wurden verschiedene
Tools der Programmiersprache Python analysiert und auf ihre Tauglichkeit im
Bezug auf TDD geprüft. Daraus entstand eine Empfehlung für vier Tools, die
gemeinsam für TDD genutzt werden können, um so optimale Funktionalität zum
Schreiben von Tests, sowie eine optimalen Testabdeckung zu garantieren. Die 
dabei erhaltenen Ergebnisse zeigen auf, dass ein Tool alleine zwar ausreichend 
für die Anwendung von TDD sein kann, jedoch ist die Verwendung von mehreren 
Tools zusammen für den Entwickler und für das Produkt besser.

Im Rahmen dieser Arbeit wurden nur jene Tools analysiert, deren letzte
Aktivität nicht weiter als ein Jahr zurück liegt. Tools, die eine Erweiterung zu
einem anderen Tool darstellen, wurden nicht behandelt, da dies den Rahmen der
Arbeit gesprengt hätte. Aus diesem Grund kann es sein, dass ein beliebtes oder
bekanntes Tool in dieser Arbeit nicht behandelt wurde. Das bedeutet nicht, dass
dieses Tool zwangsläufig nicht geeignet ist für TDD.

Im Bezug auf die Forschungsfrage, welches aktuelle Python Test-Tool das Beste 
ist, lässt sich Folgendes sagen: Viele Tools sind in ihrem Funktionsumfang sehr 
ähnlich, allerdings stellten sich \lstinline{pytest} bei den unit-testing-, 
\lstinline{mocktest} bei den mock-testing- und \lstinline{hypothesis} bei den 
fuzz-testing Tools als die beste Wahl heraus. Dennoch ergab sich, dass 
\lstinline{doctest} immer als ergänzendes Tool verwendet werden kann, egal 
welche anderen Tools verwendet werden.

Diese Arbeit hat sich jedoch lediglich mit einer begrenzten Anzahl von Tools 
beschäftigt und hat diese anhand eines spezifischen Beispiels getestet. Jeder 
Entwickler oder Software-Architekt sollte sich darüber informieren, welches 
Tool für den geplanten Anwendungszweck sinnvoll ist, da manche Tools 
vielleicht eine bessere Wahl, je nach Art der Anwendung, darstellen.
    % !TeX root = ../bachelor.tex
\section{Fazit}\label{fazit}

Durch das Testen der verschiedenen Tool der Sprache Python, hat sich diese
Arbeit mit dem Vergleich, sowie der Wertung dieser im Bezug auf TDD auseinander
gesetzt. Das Ziel dabei war es, ein oder mehrere Tool(s) zu finden, welche für
den allgemeinen Einsatz von Python und TDD am besten geeignet sind. Dieses
Ziel wurde mithilfe der gewählten Methoden erreicht und ergab, dass vier Tools
gemeinsam eine sehr stabile Basis für die Anwendung von TDD bieten. Diese vier
Tools sind: \lstinline{pytest}, \lstinline{mocktest}, \lstinline{hypothesis} und
\lstinline{doctest}.

Die dabei verwendete Testmethode ermöglicht zwar eine hohe Validität, jedoch
wäre es möglich die Reliabilität noch weiter zu steigern, indem die Testfälle
mehr Funktionen abdecken. Jedoch war die gewählte Methode für diese Arbeit die
Richtige und führt zu verwendbaren Ergebnissen.

Bei der Analyse der verschiedenen Tools ergab sich, dass die Anwendung von TDD 
komplexer ist, als zunächst erwartet. Um die Testmethode für diese Arbeit zu 
erstellen, musste eine "`Software"' programmiert werden, welche diverse 
Anforderungen erfüllen muss, da sie die Funktionen eines Tools spiegeln muss. 
Dieser erste Schritt gleicht dem ersten Schritt der Anwendung von 
TDD, da die Architektur der Software entworfen werden musste. Dies erwies sich 
schwieriger als zunächst erwartet, da die Methodik möglichst komplex sein 
sollte.

Diese Arbeit hat mir im Bezug auf TDD, Testen und Python viel Neues bei
gebracht. So würde ich zwar TDD nicht unbedingt einsetzen wollen, dies beruht
aber nicht darauf, dass ich es für ineffektiv halte, sondern vielmehr aus dem
Grund, dass es ein sehr hoher Aufwand ist, der sich erst sehr spät bezahlbar
macht. Ich selbst würde einen Ansatz nutzen, bei dem Tests neben dem 
Code geschrieben werden. Dies erspart viel Zeit und würde zu einer ähnlich 
hohen, wenn auch nicht so perfekten Testabdeckung führen.

Für mich wäre TDD nach dieser Arbeit zwar eine Methode, die ich definitiv 
einmal in einem professionellen Umfeld anwenden würde, jedoch aber eben auch 
nur dort. Meiner Meinung nach ist die Planung bei TDD das Wichtigste und sorgt 
dafür, dass die Entwickler auch Lust und Spaß an der Arbeit haben.

    \onecolumn

    % einfacher Zeilenabstand
    \singlespacing

    % Literaturverzeichnis
    % !TeX root = ../bachelor.tex
\newpage
% Literaturverzeichnis soll im Inhaltsverzeichnis auftauchen
\phantomsection
\addcontentsline{toc}{subsection}{Literaturverzeichnis}
\fancyhead[L]{Literaturverzeichnis} % Kopfzeile links
% Literaturverzeichnis endgültig anzeigen
\renewcommand\refname{Literaturverzeichnis}
\printbibliography
    
    \appendix
    \addcontentsline{toc}{section}{Anhang}
    \renewcommand{\thesubsection}{\Alph{subsection}}
    
    % das Listingverzeichnis
    % !TeX root = ../bachelor.tex
\newpage
% Listingverzeichnis soll im Inhaltsverzeichnis auftauchen
\phantomsection
\addcontentsline{toc}{section}{Listingverzeichnis}
\fancyhead[L]{Listingverzeichnis} % Kopfzeile links
\renewcommand{\lstlistlistingname}{Listingverzeichnis}
\lstlistoflistings
    
    % die Abbildungen
    % das Abbildungsverzeichnis
\newpage
% Abbildungsverzeichnis soll im Inhaltsverzeichnis auftauchen
\phantomsection
\addcontentsline{toc}{subsection}{Abbildungsverzeichnis}
\fancyhead[L]{Abbildungsverzeichnis} % Kopfzeile links
% Abbildungsverzeichnis endgültig anzeigen
\listoffigures
    
    % das Fußnotenverzeichnis
    % !TeX root = ../bachelor.tex
\newpage
% Fußnotenverzeichnis soll im Inhanltsverzeichnis auftauchen
\phantomsection
\addcontentsline{toc}{subsection}{Fußnotenverzeichnis}
\fancyhead[L]{Fußnotenverzeichnis} % Kopfzeile links
\fancyhead[L]{\nouppercase{\leftmark}} % Kopfzeile links
% Fußnotenverzeichnis endgültig anzeigen
\listoffootnotes
    
    % das Glossar
    \newpage
    % Glossar soll im Inhaltsverzeichnis auftauchen
    \phantomsection
    \addcontentsline{toc}{subsection}{Glossar}
    \fancyhead[L]{Glossar} % Kopfzeile links
    % das Abkürzungsverzeichnis entgültige Ausgeben
    \printglossaries
    
    % das Abkürzungsverzeichnis
    % !TeX root = ../../main.tex
\nomenclature{GPL}{GNU General Public License}
\nomenclature{GNU}{GNU is not Unix}
\nomenclature{LGPL}{GNU Lesser General Public License}
\nomenclature{TDD}{Test-driven development}
\nomenclature{GUI}{Graphical User Interface}
\nomenclature{STDLIB}{Python Standard Bibliothek}
\nomenclature{VCS}{\gls{vcs}}
    
    % die Listings
    % !TeX root = ../bachelor.tex
\newpage
% Die Listings sollen im Inhanltsverzeichnis auftauchen
\section*{Listings}
\addcontentsline{toc}{section}{Listings}
\fancyhead[L]{Listings} % Kopfzeile links

%=====================================================================%
% Base Modules
%=====================================================================%
\lstinputlisting[
    language=Python,
    label=listing:base:my_module,
    caption=Basis Modul zum testen,
    lastline=42
]{analyse/my_package/my_module.py}

\lstinputlisting[
    language=Python,
    label=listing:base:my_mock_module,
    caption=Basis Modul zum testen von mocks,
]{analyse/my_package/my_mock_module.py}
\newpage

%=====================================================================%
% unittest
%=====================================================================%
\lstinputlisting[
    language=Python,
    label=listing:unittest:example,
    caption=unittest einfaches Beispiel
]{analyse/unittest_/example.py}

\lstinputlisting[
    label=listing:unittest:example_output_success,
    caption=unittest einfaches Beispiel: Output erfolgreich
]{analyse/unittest_/example_output_success.txt}

\lstinputlisting[
    label=listing:unittest:example_output_failure,
    caption=unittest einfaches Beispiel: Output misslungen
]{analyse/unittest_/example_output_failure.txt}
\newpage

\lstinputlisting[
    language=Python,
    label=listing:unittest:advanced,
    caption=unittest my\_module
]{analyse/unittest_/advanced_example.py}
\newpage

%=====================================================================%
% doctest
%=====================================================================%
\lstinputlisting[
    language=Python,
    label=listing:doctest:example,
    caption=doctest unterscheidet Typen
]{analyse/doctest/example.py}

\lstinputlisting[
    label=listing:doctest:example_output,
    caption=doctest unterscheidet Typen: Output
]{analyse/doctest/example_output.txt}

\lstinputlisting[
    label=listing:doctest:output,
    caption=doctest verbose Output
]{analyse/doctest/output.txt}
\newpage

\lstinputlisting[
    language=Python,
    label=listing:doctest:advanced,
    caption=doctest: my\_module,
    firstline=19
]{analyse/doctest/advanced.py}

\lstinputlisting[
    language=Python,
    label=listing:doctest:advanced_text,
    caption=doctest my\_module: Textdatei
]{analyse/doctest/advanced.txt}
\newpage

%=====================================================================%
% pytest
%=====================================================================%
\lstinputlisting[
    language=Python,
    label=listing:pytest:yield,
    caption=pytest setUp und tearDown
]{analyse/py.test/yield.py}

\lstinputlisting[
    label=listing:pytest:yield_output,
    caption=pytest setUp und tearDown Output
]{analyse/py.test/yield.txt}

\lstinputlisting[
    language=Python,
    label=listing:pytest:advanced,
    caption=pytest my\_module
]{analyse/py.test/advanced.py}

\lstinputlisting[
    label=listing:pytest:advanced_output,
    caption=pytest my\_module: Output
]{analyse/py.test/advanced_output.txt}

\lstinputlisting[
label=listing:pytest:requirements,
caption=pytest Abhängigkeiten
]{analyse/py.test/after.requirements}

%=====================================================================%
% reahl.stubble
%=====================================================================%
\lstinputlisting[
    language=Python,
    label=listing:reahl.stubble:Impostor,
    caption=stubble: Impostor
]{analyse/stubble/impostor.py}

\lstinputlisting[
    language=Python,
    label=listing:reahl.stubble:example,
    caption=stubble my\_mock\_module
    lastline=91
]{analyse/stubble/example.py}

%=====================================================================%
% mocktest
%=====================================================================%
\lstinputlisting[
    language=Python,
    label=listing:mocktest:example,
    caption=mocktest my\_mock\_module
]{analyse/mocktest_/example.py}

%=====================================================================%
% flexmock
%=====================================================================%
\lstinputlisting[
    language=Python,
    label=listing:flexmock:example,
    caption=flexmock my\_mock\_module
]{analyse/flexmock_/example.py}

\lstinputlisting[
    label=listing:flexmock:exception,
    caption=flexmock exception
]{analyse/flexmock_/exception.txt}

%=====================================================================%
% doublex
%=====================================================================%
\lstinputlisting[
    language=Python,
    label=listing:doublex:example,
    caption=python-doublex my\_mock\_module
]{analyse/python-doublex/example.py}

%=====================================================================%
% hypothesis
%=====================================================================%
\lstinputlisting[
    language=Python,
    label=listing:hypothesis:given,
    caption=hypothesis @given() Beispiel
]{analyse/hypothesis_/given.py}
    
    % 1,5 facher Zeilenabstand
    \onehalfspacing
    
    % Eidesstattliche Erklärung
    % !TeX root = bachelor.tex
\newpage

\thispagestyle{empty}
\addcontentsline{toc}{section}{Eidesstattliche Erklärung}

\begin{verbatim}

\end{verbatim}

\begin{LARGE}
\begin{center}
	\textbf{Eidesstattliche Erklärung zur Abschlussarbeit}
\end{center}
\end{LARGE}
\begin{verbatim}


\end{verbatim}
Hiermit versichere ich, die eingereichte Abschlussarbeit selbständig verfasst und
keine andere als die von mir angegebenen Quellen und Hilfsmittel benutzt zu haben.
Wörtlich oder inhaltlich verwendete Quellen wurden entsprechend den anerkannten
Regeln wissenschaftlichen Arbeitens zitiert. Ich erkläre weiterhin, dass die
vorliegende Arbeit noch nicht anderweitig als Abschlussarbeit eingereicht wurde.
\\
\\
Das Merkblatt zum Täuschungsverbot im Prüfungsverfahren der Hochschule
Augsburg habe ich gelesen und zur Kenntnis genommen. Ich versichere, dass die
von mir abgegebene Arbeit keinerlei Plagiate, Texte oder Bilder umfasst, die durch
von mir beauftragte Dritte erstellt wurden.
\begin{verbatim}



\end{verbatim}


\begin{displaymath}
% use packages: array
\begin{array}{ll}
Unterschrift:~~~~~~~~~~~~~~~~~~~~~~~~~~~~~~~~~~~~~~~~~~
& Ort, Datum:~~~~~~~~~~~~~~~~~~~~~~~~~~~~~~~~~~~~~~~~~~
\end{array}
\end{displaymath}
\end{document}

% Indexerstellung
\makeindex
